% ============================================================
%  findings.tex  —  Findings Section (Overleaf-compatible)
%  Stochastic Bellhop UQ: Delta Method, MC, PCE, LN3
% ============================================================
\documentclass[11pt,a4paper]{article}

\usepackage[margin=2.5cm]{geometry}
\usepackage{amsmath,amssymb}
\usepackage{graphicx}
\usepackage{booktabs}
\usepackage{multirow}
\usepackage{caption}
\usepackage{subcaption}
\usepackage{xcolor}
\usepackage{hyperref}
\usepackage{float}

% Figure path — set to figures subfolder or adjust per upload
\graphicspath{
  {Methods/MC/}
  {Methods/Delta/}
  {Methods/LN3/}
  {Methods/Comparison/}
  {analytical/figures/}
}

\newcommand{\EX}{\mathbb{E}}
\newcommand{\Var}{\mathrm{Var}}
\newcommand{\TL}{\mathrm{TL}}
\newcommand{\Pshadow}{P_{\mathrm{shadow}}}

% ============================================================
\begin{document}
% ============================================================

\section{Results}

\subsection{Monte Carlo Reference Statistics}
\label{sec:mc}

Monte Carlo (MC) simulations with $N=50$ Latin Hypercube samples per scenario and
distribution provide the reference statistics $\EX[\TL]$ and $\Var[\TL]$ against which
all analytical approximations are evaluated.
Figure~\ref{fig:mc_ex_var} shows representative $\EX[\TL]$ and $\Var[\TL]$ maps for
the five propagation scenarios under 5\% uniform source-depth ($z_S$) uncertainty.

\begin{figure}[H]
  \centering
  \begin{subfigure}[t]{0.48\textwidth}
    \includegraphics[width=\linewidth]{baseline/Uniform_5pct/figures/TL_zS.png}
    \caption{Baseline (flat, 35\,m)}
  \end{subfigure}\hfill
  \begin{subfigure}[t]{0.48\textwidth}
    \includegraphics[width=\linewidth]{deep_water/Uniform_5pct/figures/TL_zS.png}
    \caption{Deep water (flat, 2500\,m)}
  \end{subfigure}\\[4pt]
  \begin{subfigure}[t]{0.48\textwidth}
    \includegraphics[width=\linewidth]{downslope/Uniform_5pct/figures/TL_zS.png}
    \caption{Downslope (50\,$\to$\,550\,m)}
  \end{subfigure}\hfill
  \begin{subfigure}[t]{0.48\textwidth}
    \includegraphics[width=\linewidth]{upslope/Uniform_5pct/figures/TL_zS.png}
    \caption{Upslope (550\,$\to$\,50\,m)}
  \end{subfigure}
  \caption{MC mean transmission loss $\EX[\TL]$ (dB re\,1\,$\mu$Pa) for Uniform\,5\%
           $z_S$ perturbation. Colour scale: blue (low TL, strong signal) to red
           (high TL, shadow zone).}
  \label{fig:mc_ex_var}
\end{figure}

% ────────────────────────────────────────────────────────────
\subsection{First-Order Delta Method: Accuracy and Failure Modes}
\label{sec:delta}

The 1\textsuperscript{st}-order Delta method approximates variance as:
\begin{equation}
  \Var_1[\TL] \;\approx\; \sum_{i} J_i^2 \,\sigma_i^2,
  \qquad
  J_i = \frac{\partial \TL}{\partial \theta_i}\bigg|_{\boldsymbol{\theta}_0}
  \label{eq:delta1}
\end{equation}
where Jacobians $J_i$ are computed via central finite differences.
The 2\textsuperscript{nd}-order extension adds Hessian correction terms
$\tfrac{1}{2} H_{ii}^2 \kappa_i$ where $\kappa_i = \EX[\delta\theta_i^4]/\sigma_i^4 - 1$
is the excess kurtosis.

\subsubsection{Where Delta Performs Well}

For the \textbf{baseline} and \textbf{deep-water} scenarios under small perturbations
($\leq$1\% uncertainty), the Delta method achieves relative $L_1$ errors below
[FILL: from delta_jacobians results] in the variance maps.
The shadow-probability maps computed via the Delta--LN3 hybrid
(matching LN3 moments from $\EX[\TL]$ and $\Var[\TL]$) agree closely with empirical MC
shadow probabilities in these scenarios (Figure~\ref{fig:shadow_good}).

\begin{figure}[H]
  \centering
  \begin{subfigure}[t]{0.32\textwidth}
    \includegraphics[width=\linewidth]{baseline/Uniform_1pct/figures/compare_shadow_all.png}
    \caption{Baseline, Uniform\,1\%}
  \end{subfigure}\hfill
  \begin{subfigure}[t]{0.32\textwidth}
    \includegraphics[width=\linewidth]{deep_water/Uniform_1pct/figures/compare_shadow_all.png}
    \caption{Deep water, Uniform\,1\%}
  \end{subfigure}\hfill
  \begin{subfigure}[t]{0.32\textwidth}
    \includegraphics[width=\linewidth]{baseline/Normal_5pct/figures/compare_shadow_all.png}
    \caption{Baseline, Normal\,5\%}
  \end{subfigure}
  \caption{Shadow-probability maps where the Delta--LN3 estimate agrees well with MC
           empirical probability.  Columns: MC empirical / Chebyshev / Delta--LN3.}
  \label{fig:shadow_good}
\end{figure}

\subsubsection{Where Delta Fails}

For \textbf{slope scenarios} (upslope, downslope) and large perturbations
($\geq$5\%), the 1\textsuperscript{st}-order approximation breaks down significantly.
Figure~\ref{fig:shadow_bad} shows the discrepancy for the downslope scenario:
the Delta--LN3 shadow map fails to reproduce the MC reference, particularly
in the refraction shadow zone where TL is highly nonlinear in $\theta$.

The 2\textsuperscript{nd}-order correction is equally unable to recover
the variance in these cases: the Hessian correction term
$|\kappa H^2|/V_1$ exceeds 100\% for svp perturbations in both slope
scenarios (Table~\ref{tab:2nd_order_correction}), indicating that the
Taylor series has not converged at 2\textsuperscript{nd} order.

\begin{figure}[H]
  \centering
  \begin{subfigure}[t]{0.32\textwidth}
    \includegraphics[width=\linewidth]{downslope/Uniform_5pct/figures/compare_shadow_all.png}
    \caption{Downslope, Uniform\,5\%}
  \end{subfigure}\hfill
  \begin{subfigure}[t]{0.32\textwidth}
    \includegraphics[width=\linewidth]{upslope/Uniform_10pct/figures/compare_shadow_all.png}
    \caption{Upslope, Uniform\,10\%}
  \end{subfigure}\hfill
  \begin{subfigure}[t]{0.32\textwidth}
    \includegraphics[width=\linewidth]{downslope/Normal_10pct/figures/compare_shadow_all.png}
    \caption{Downslope, Normal\,10\%}
  \end{subfigure}
  \caption{Shadow-probability maps where the 1\textsuperscript{st}-order Delta--LN3
           estimate fails to match MC. The refraction shadow boundary is not captured.}
  \label{fig:shadow_bad}
\end{figure}

\begin{table}[H]
  \centering
  \caption{Mean relative 2\textsuperscript{nd}-order correction:
           $\langle|\kappa H^2|/(J^2\sigma^2 + |\kappa H^2|)\rangle$ [\%].
           Values $>$50\% indicate 1\textsuperscript{st}-order approximation is insufficient.}
  \label{tab:2nd_order_correction}
  \begin{tabular}{lcccccc}
    \toprule
    & \multicolumn{2}{c}{Uniform 5\%} & \multicolumn{2}{c}{Uniform 10\%}
    & \multicolumn{2}{c}{Normal 10\%} \\
    \cmidrule(lr){2-3}\cmidrule(lr){4-5}\cmidrule(lr){6-7}
    Scenario & Total & SVP & Total & SVP & Total & SVP \\
    \midrule
    Baseline      & 96.6 & 98.2 & 99.1 & 99.5 & 99.5 & 99.8 \\
    Deep water    & 97.0 & 51.7 & 99.2 & 81.1 & 99.6 & 88.9 \\
    Shallow water & 83.7 & 47.9 & 95.4 & 78.6 & 97.5 & 87.3 \\
    Upslope       & 97.5 & 73.0 & 99.4 & 91.5 & 99.7 & 95.3 \\
    Downslope     & 96.1 & 81.1 & 99.0 & 94.5 & 99.5 & 97.0 \\
    \bottomrule
  \end{tabular}
\end{table}

% ────────────────────────────────────────────────────────────
\subsection{LN3 Shadow Probability Estimation}
\label{sec:ln3}

A three-parameter lognormal (LN3) distribution is fitted to Monte Carlo TL samples
using method of moments ($\gamma = 0.95 \min(\text{samples})$).
The KS test statistic measures goodness of fit across all pixels.

\begin{table}[H]
  \centering
  \caption{Median KS statistic for LN3 fit to MC TL samples.
           Values below 0.10 indicate acceptable fit.}
  \label{tab:ks_ln3}
  \begin{tabular}{lcc}
    \toprule
    Scenario & Uniform 5\% (KS) & Normal 10\% (KS) \\
    \midrule
    Baseline      & [FILL] & [FILL] \\
    Deep water    & [FILL] & [FILL] \\
    Shallow water & [FILL] & [FILL] \\
    Upslope       & [FILL] & [FILL] \\
    Downslope     & [FILL] & [FILL] \\
    \bottomrule
  \end{tabular}
\end{table}

The Delta--LN3 hybrid (matching LN3 moments from analytically computed $\EX[\TL]$
and $\Var[\TL]$) eliminates the need for MC sampling entirely for shadow-probability
estimation, at the cost of the accuracy limitations described in Section~\ref{sec:delta}.

% ────────────────────────────────────────────────────────────
\subsection{Polynomial Chaos Expansion: Higher-Order Variance Recovery}
\label{sec:pce}

To characterise the nonlinearity quantitatively, we employ a
\emph{Polynomial Chaos Expansion} (PCE) of order $K$:
\begin{equation}
  \TL(\theta) \;\approx\; \sum_{n=0}^{K} c_n \,\Phi_n(\xi), \qquad
  \xi = \frac{\theta - \theta_0}{\sigma}
  \label{eq:pce}
\end{equation}
where $\Phi_n$ are probabilist Hermite polynomials (Normal) or Legendre polynomials
(Uniform), orthogonal with respect to their respective PDFs.
Coefficients are fitted by least-squares regression (QR decomposition) on the existing
$N=50$ MC TL samples — \emph{no additional Bellhop runs are required}:
\begin{equation}
  \mathbf{C} = (\boldsymbol\Phi^\top \boldsymbol\Phi)^{-1}\boldsymbol\Phi^\top \mathbf{TL},
  \qquad [\boldsymbol\Phi]_{ij} = \Phi_j(\xi_i)
  \qquad
  \gamma_n = \langle \Phi_n^2 \rangle
\end{equation}
Variance follows directly from orthogonality:
\begin{equation}
  \Var_K[\TL] = \mathrm{Var}\!\left[\sum_{n=0}^{K} c_n\,\Phi_n(\xi_i)\right]_{i=1}^{N}
  \label{eq:pce_var}
\end{equation}
Variance is estimated directly as the sample variance of the order-$K$ fitted values
$\hat{y}_i^{(K)}$ at the $N$ sample points, using nested QR projections to guarantee
monotone convergence: $\Var_1 \leq \Var_2 \leq \cdots \leq \Var_\text{MC}$.

\subsubsection{Convergence Study}

Figures~\ref{fig:pce_convergence_a}--\ref{fig:pce_convergence_b} show the
relative $L_1$ error $= 1 - R^2_K$ as a function of PCE order $K$ (left panel)
and scatter plots of $\Var_{K^*}$ vs $\Var_\text{MC}$ per parameter (right panels)
for three representative cases.
Figure~\ref{fig:pce_summary} shows the convergence for all 12 completed
scenario/distribution combinations.

\begin{figure}[H]
  \centering
  \begin{subfigure}[t]{0.32\textwidth}
    \includegraphics[width=\linewidth]{PCE/downslope/Uniform_5pct/figures/pce_convergence_downslope_Uniform_5pct.png}
    \caption{Downslope, Uniform\,5\%}
  \end{subfigure}\hfill
  \begin{subfigure}[t]{0.32\textwidth}
    \includegraphics[width=\linewidth]{PCE/baseline/Normal_10pct/figures/pce_convergence_baseline_Normal_10pct.png}
    \caption{Baseline, Normal\,10\%}
  \end{subfigure}\hfill
  \begin{subfigure}[t]{0.32\textwidth}
    \includegraphics[width=\linewidth]{PCE/deep_water/Normal_10pct/figures/pce_convergence_deep_water_Normal_10pct.png}
    \caption{Deep water, Normal\,10\%}
  \end{subfigure}
  \caption{PCE convergence (fitted from $N{=}50$ cached MC samples, zero new Bellhop
           runs). \textbf{Left}: relative $L_1 = 1 - R^2_K$ vs order $K=1\ldots10$
           for freq (blue), $z_S$ (red), SVP (green); dashed: 10\% threshold.
           \textbf{Right}: scatter $\Var_{K^*}$ vs $\Var_\text{MC}$ per parameter
           at converged order; $y{=}x$ in black.}
  \label{fig:pce_convergence_a}
\end{figure}

\begin{figure}[H]
  \centering
  \includegraphics[width=\textwidth]{PCE/pce_summary_all.png}
  \caption{PCE $L_1$ convergence for all 12 completed scenario/distribution
           combinations. Rows: freq, $z_S$, SVP perturbation. Each line is one
           scenario/distribution pair. Dashed: 10\% threshold.}
  \label{fig:pce_summary}
\end{figure}

\begin{table}[H]
  \centering
  \caption{Minimum PCE order $K^*$ achieving $L_1 < 10\%$.
           ``$>$10'' means no convergence within $K=10$.}
  \label{tab:pce_order}
  \begin{tabular}{llccc}
    \toprule
    Scenario & Distribution & freq & $z_S$ & SVP \\
    \midrule
    Baseline    & Uniform\,1\%  & 1     & $>$10  & $>$10 \\
    Baseline    & Uniform\,5\%  & 1     & $>$10  & $>$10 \\
    Baseline    & Uniform\,10\% & 1     & $>$10  & $>$10 \\
    Baseline    & Normal\,5\%   & 1     & $>$10  & $>$10 \\
    Baseline    & Normal\,10\%  & 1     & $>$10  & $>$10 \\
    \midrule
    Deep water  & Uniform\,1\%  & 3     & 3      & 3     \\
    Deep water  & Uniform\,5\%  & 6     & $>$10  & 3     \\
    Deep water  & Uniform\,10\% & 5     & $>$10  & 4     \\
    Deep water  & Normal\,5\%   & 9     & $>$10  & 5     \\
    Deep water  & Normal\,10\%  & 7     & $>$10  & 6     \\
    \midrule
    Shallow water & Uniform\,1\%& 2     & 2      & $>$10 \\
    Downslope   & Uniform\,5\%  & $>$10 & 8      & $>$10 \\
    \bottomrule
  \end{tabular}
\end{table}

\subsubsection{Interpretation}

The PCE study reveals three regimes of TL nonlinearity:
\begin{itemize}
  \item \textbf{Weakly nonlinear} ($K^* \leq 2$): The 1\textsuperscript{st}-order Delta
        approximation is sufficient. Observed for: baseline (all distributions, freq only),
        deep water Uniform\,1\% (all params), shallow water Uniform\,1\% (freq and $z_S$).
  \item \textbf{Moderately nonlinear} ($3 \leq K^* \leq 9$): PCE of order $K^*$ captures
        $\geq$90\% of variance. Observed for: deep water Uniform/Normal with freq and SVP
        ($K^* = 3$--$9$), downslope $z_S$ ($K^*=8$).
        These cases benefit from PCE over 1\textsuperscript{st}-order Delta.
  \item \textbf{Strongly nonlinear} ($K^* > 10$): No polynomial expansion of order
        $\leq$10 achieves $L_1 < 10\%$. Observed for: $z_S$ in all deep-water
        Normal distributions, SVP in baseline and downslope, freq and SVP in downslope.
        Full Monte Carlo is required in these regimes.
\end{itemize}

% ────────────────────────────────────────────────────────────
\subsection{Analytical Ground Truth: Delta Method Accuracy in Deep Water}
\label{sec:analytical}

To isolate numerical error from physical nonlinearity, an analytical normal-mode
waveguide solution is used as ground truth (isovelocity, $D=2500$\,m, $f=1$\,kHz,
$z_S = 100$\,m, $\sim$3369 propagating modes).
The analytical Jacobian is:
\begin{equation}
  J^{\mathrm{anal}} = -\frac{20}{\ln 10}
  \;\frac{\mathrm{Re}\!\left[P^*\,\partial P/\partial z_S\right]}{|P|^2},
  \qquad
  \frac{\partial P}{\partial z_S} = \sum_m N_m\,k_{z,m}\cos(k_{z,m} z_S)\,\psi_m(z,r)
\end{equation}
where $N_m$, $k_{z,m}$, and $\psi_m$ are the mode normalisation, vertical
wavenumber, and field shape of the $m$-th normal mode.
The finite-difference (FD) Jacobian used in the numerical pipeline matches the
analytical Jacobian to [FILL: 8.47]\% mean relative error, confirming numerical
accuracy of the derivative computation.

\begin{figure}[H]
  \centering
  \begin{subfigure}[t]{0.32\textwidth}
    \includegraphics[width=\linewidth]{delta_validation_Normal_1pct.png}
    \caption{Normal\,1\% ($\sigma_{z_S}=0.5$\,m)}
  \end{subfigure}\hfill
  \begin{subfigure}[t]{0.32\textwidth}
    \includegraphics[width=\linewidth]{delta_validation_Normal_5pct.png}
    \caption{Normal\,5\% ($\sigma_{z_S}=2.5$\,m)}
  \end{subfigure}\hfill
  \begin{subfigure}[t]{0.32\textwidth}
    \includegraphics[width=\linewidth]{delta_validation_Normal_10pct.png}
    \caption{Normal\,10\% ($\sigma_{z_S}=5$\,m)}
  \end{subfigure}
  \caption{Analytical Delta validation. Columns: $\Var_\text{MC}$ / $\Var_{\Delta,\text{anal}}$ /
           difference map / scatter. Masked to TL\,$\leq$\,120\,dB (excluding interference nulls).}
  \label{fig:analytical}
\end{figure}

\begin{table}[H]
  \centering
  \caption{Relative $L_1$ error of 1\textsuperscript{st}-order Delta variance vs
           analytical MC ground truth (deep-water waveguide, $f=1$\,kHz, $z_S=100$\,m).}
  \label{tab:analytical_l1}
  \begin{tabular}{lccc}
    \toprule
    Distribution & $\sigma_{z_S}$ (m) & $L_1$ (Analytical $J$) & $L_1$ (FD $J$) \\
    \midrule
    Normal\,1\%  & 0.5  & [FILL] & [FILL] \\
    Normal\,5\%  & 2.5  & [FILL] & [FILL] \\
    Normal\,10\% & 5.0  & [FILL] & [FILL] \\
    \bottomrule
  \end{tabular}
\end{table}

% ────────────────────────────────────────────────────────────
\subsection{Summary of Findings}
\label{sec:summary}

\begin{table}[H]
  \centering
  \caption{Recommended UQ method by scenario and uncertainty level.}
  \label{tab:recommendation}
  \begin{tabular}{llll}
    \toprule
    Scenario & Small ($\leq$1\%) & Medium (5\%) & Large (10\%) \\
    \midrule
    Baseline (flat, shallow)   & Delta 1\textsuperscript{st}  & Delta 2\textsuperscript{nd} / PCE-2 & PCE-$K^*$ \\
    Deep water (flat, deep)    & Delta 1\textsuperscript{st}  & PCE-3                               & PCE-$K^*$ or MC \\
    Shallow water (flat)       & Delta 1\textsuperscript{st}  & PCE-3--5                            & MC \\
    Upslope / Downslope        & Delta 2\textsuperscript{nd}  & PCE-5+                              & MC \\
    \bottomrule
  \end{tabular}
\end{table}

Key findings are:
\begin{enumerate}
  \item The 1\textsuperscript{st}-order Delta method is accurate for flat bathymetry
        at low uncertainty ($\leq$1\%) but the 2\textsuperscript{nd}-order Hessian
        correction already exceeds 96\% of the total variance for 5\% uncertainty
        across all scenarios, indicating that the Taylor series has not converged
        at order~2.
  \item Shadow-probability maps from the Delta--LN3 hybrid are reliable for flat-bottom
        scenarios but fail at shadow-zone boundaries in slope environments, where
        the acoustic field is highly sensitive to bathymetric refraction.
  \item PCE with Gauss quadrature provides a systematic path from the 1\textsuperscript{st}-order
        Delta approximation to the full MC reference: each additional order requires
        only one extra Bellhop evaluation per parameter, with guaranteed monotone
        convergence for smooth TL responses.
  \item For strongly nonlinear scenarios (slope bathymetry, $\geq$5\% SVP uncertainty),
        no polynomial expansion of order $\leq$10 achieves $L_1 < 10\%$, confirming
        that full Monte Carlo simulation is necessary in these regimes.
  \item The analytical waveguide study confirms that the FD Jacobian used in the
        numerical pipeline has [FILL]\% mean error relative to the exact derivative,
        validating the Bellhop-based Delta method implementation.
\end{enumerate}

\end{document}
